\rok{Децембар 2025.}{E}

Први практични тест из Алгоритама и структура података

\zadatak{1}{Heap Sort}
Написати Heap Sort алгоритам за сортирање целобројног низа дужине $n$.

\textbf{Додатне напомене за решавање задатка:}
\begin{itemize}
	\item Улаз је несортиран низ целих бројева.
	\item Излаз је сортирани низ целих бројева.
	\item За имплементацију алгоритма користити класу \texttt{Heap} која се налази у template-у.
	\item У template-у се налази класа \texttt{RandomArrayGenerator} коју треба искористити за генерисање низа случајних бројева.
	\item Креирати класу \texttt{HeapSort} и у оквиру ње генерисати методу:
	\begin{Verbatim}[fontsize=\small]
public static void sort(long[] arr)
	\end{Verbatim}
\end{itemize}

\zadatak{2}{Кумулативна сума вредности левих суседа}
Написати алгоритам који модификује листу тако што вредност чвора замењује вредношћу суме свих елемената који се у листи налазе са његове леве стране.

\textbf{Додатне напомене за решавање задатка:}
\begin{itemize}
	\item Захтеван потпис методе: \texttt{public void calculateCumulativeSum()}.
	\item Методу писати у оквиру класе \texttt{LinkedList}.
\end{itemize}

\textbf{Потпис методе:}
\begin{Verbatim}[fontsize=\small]
public void calculateCumulativeSum()
\end{Verbatim}

\textbf{Пример:}
\begin{itemize}
	\item \textbf{Улаз:} \texttt{LinkedList: 1 -> 0 -> 6 -> 8 -> 4} \\
	      \textbf{Излаз:} \texttt{0 -> 1 -> 1 -> 7 -> 15}
	\item \textbf{Улаз:} \texttt{LinkedList: 1 -> 2 -> 4 -> 0 -> 6} \\
	      \textbf{Излаз:} \texttt{0 -> 1 -> 3 -> 7 -> 7}
\end{itemize}

\zadatak{3}{Hot Potato Game}
\textbf{Проблем:} У класичној игри "Врући кромпир", играчи у кругу пребацују кромпир док музика свира. Када музика стане, играч који држи кромпир испада. У нашој напредној варијанти, број пролазака кромпира се динамички мења током игре.

\textbf{Правила игре:}
\begin{enumerate}
	\item Има $n$ играча нумерисаних од 1 до $n$.
	\item Почетни број пролазака је $K$.
	\item Након сваког испадања играча, број пролазака за следећи круг се мења по формули:
	$$
	\text{нови\_проласци} = (\text{стари\_проласци} \times 2) \% (\text{преостали\_играчи} + 1)
	$$
	\item Ако је резултат формуле 0, тада је \texttt{novi\_prolasci = preostali\_igraci}.
	\item Играчи се елиминишу док не остане само један - победник.
\end{enumerate}

\textbf{Додатни захтеви за решавање задатка:}
\begin{enumerate}
	\item Задатак решавати у оквиру методе:
	\begin{Verbatim}[fontsize=\small]
public int PlayHotPotato(int totalPlayers, int initialPasses)
	\end{Verbatim}
	\item Морате користити Queue за представљање круга играча.
	\item Имплементирати динамичко рачунање пролазака по датој формули.
	\item Резултат извршавања алгоритма треба да буде редни број победника у игри.
\end{enumerate}

\textbf{Пример:}
\begin{itemize}
	\item \textbf{Улаз 1:} \texttt{totalPlayers = 7, initialPasses = 3} \\
	      \textbf{Излаз 1:} \texttt{expectedWinner = 4}
	\item \textbf{Улаз 2:} \texttt{totalPlayers = 10, initialPasses = 1} \\
	      \textbf{Излаз 2:} \texttt{expectedWinner = 8}
\end{itemize}

\sablon{Шаблон првог теста децембар 2025. група E}{2025/Decembar/GrupaE/AISP2025_Test1_GrupaE.linq}
